\clearpage
\section{Secular terms in weakly non-linear oscillators}
the equation is 
\begin{equation}
    \ddot{x}+x = -\epsilon\ins{a_2 x^2+a_3x^3},\quad 0<\epsilon\ll1
\end{equation}
\subsection{two time expansion}
we define 
\begin{equation}
    \tau = t \quad T=\epsilon t
\end{equation}
so that 
\begin{equation}
    \del{t}=\del{\tau}+\epsilon\del{T}
\end{equation}
and 
\begin{equation}
    \del{t}^2=\del{\tau}^2+2\epsilon\del{\tau}\del{T}+O\ins{\epsilon^2}
\end{equation}
we expand 
\begin{equation}
    x = x_0(\tau,T)+\epsilon x_1(\tau, T)+O(\epsilon^2)
\end{equation}
substituting into the equation
\begin{equation}
    \ins{\del{\tau}^2+2\epsilon\del{\tau}\del{T}}\ins{x_0+\epsilon x_1}+x_0+\epsilon x_1 = -\epsilon\ins{a_2x_0^2+a_3x_0^3}+O(\epsilon^2)
\end{equation}
and at \(O(1)\)
\begin{equation}
    \del{\tau}^2 x_0+x_0=0
\end{equation}
whose solution can be written as 
\begin{equation}
    x_0\ins{\tau,T}=r(T)\cos\psi,\quad \psi=\tau+\phi(T)
\end{equation}
where \(r\) is the amplitude and \(\phi\) is the phase. At order \(O(\epsilon)\)
\begin{equation}
    \del{\tau}^2x_1+x_1 = -2\del{\tau}\del{T}x_0-a_2x_0^2-a_3x_0^3
\end{equation}
\subsection{higher harmonics}
Starting from 
\begin{equation}
    x_0 = r\cos\psi
\end{equation}
then 
\begin{equation}
    \del{T}x_0 = \del{T}r\cos\psi-r\del{T}\phi\sin\psi
\end{equation}
then 
\begin{equation}
    \del{\tau}\del{T}x_0 = -\del{T}r\sin\psi-r\del{T}\phi\cos\psi
\end{equation}
so 
\begin{equation}
    -2\del{\tau}\del{T}x_0 = \del{T}r\sin\psi+r\del{T}\phi\cos\psi
\end{equation}
and using the identities
\begin{align}
    \cos^2\psi=\frac{1}{2}\ins{1+\cos2\psi} && \cos^3\psi=\frac{1}{4}\ins{3\cos\psi+\cos3\psi}
\end{align}
we get 
\begin{equation}
    -a_2x_0^2=-\frac{a_2r^2}{2}-\frac{a_2r^2}{2}\cos2\psi
\end{equation}
and 
\begin{equation}
    -a_3x_0^3=-\frac{3a_3r^3}{4}\cos\psi-\frac{a_3r^3}{4}\cos3\psi
\end{equation}
combining everything
\begin{equation}
    \del{\tau}^2x_1+x_1 = \del{T}r\sin\psi+\ins{2r\del{T}\phi-\frac{3a_3r^3}{4}}\cos\psi-\frac{a_2r^2}{2}\ins{1+\cos2\psi}-\frac{a_3r^3}{4}\cos3\psi
\end{equation}
\subsection{resonant terms}
The resonant terms are \(2\del{T}r\sin\psi\) and \(\ins{2r\del{T}\phi-\frac{3a_3e^3}{4}}\cos\psi\). To eliminate secular growth their coefficients must vanish so 
\begin{equation}
    \del{T}r=0 \implies r(T)=r_0
\end{equation}
and
\begin{equation}
    \del{T}\phi=\frac{3a_3r^2}{8} \implies \phi(T)=\phi_0+\frac{3}{8}a_3r_0^2T
\end{equation}
inserting \(T=\epsilon t\)
\begin{equation}
    x(t) \simeq r_0\cos\ins{\ins{1+\frac{3}{8}\epsilon a_3 r_0^2}t+\phi_0}
\end{equation}
\subsection{duffing oscillator}
in the duffing oscillator
\begin{equation}
    a_2=0,\quad a_3=1
\end{equation}
giving
\begin{equation}
    \del{T}r=0,\quad \del{T}\phi=\frac{3}{8}r^2
\end{equation}
A nonzero quadratic nonlinearity \(a_2\) doesn't change the leading order slow equations because 
\begin{equation}
    x_0^2=\frac{r^2}{2}\ins{1+\cos2\psi}
\end{equation}
only has a constant term and a second harmonic, and since it doesn't have the natural frequency of the oscillator, it's not resonant.
\subsection{averaged equations}
The averaged equations from the notes are 
\begin{equation}
    \dot{A}=\epsilon\left<h\sin\theta\right>,\quad A\dot{\phi}=\epsilon\left<h\cos\theta\right>, \quad \left<\cdot\right>=\frac{1}{2\pi}\int_{-\pi}^{\pi}\ins{\cdot}\dd\theta
\end{equation}
inserting 
\begin{equation}
    x=A\cos\theta,\quad\dot{x}=-A\sin\theta
\end{equation}
where \(\theta=t+\phi(t)\) into \(h\)
\begin{equation}
    h = A^3\cos\theta\sin^2\theta
\end{equation}
making the amplitude equation
\begin{equation}
    \dot{A}=\epsilon A^3 \left<\cos\theta\sin^3\theta\right>
\end{equation}
but the average over a period of the trig function zeroes
\begin{equation}
    \dot{A}=0
\end{equation}
and for the phase 
\begin{align}
    A\dot{\phi}&=\epsilon A^3\left<\cos^2\theta\sin^2\theta\right>\\
    &= \epsilon A^3\left<\frac{1}{4}\sin^2 2\theta\right>\\
    &= \frac{\epsilon A^3}{4}\cdot\frac{1}{2}\\
    &= \frac{\epsilon A^3}{8}
\end{align}
therefore 
\begin{equation}
    \dot{\phi} = \frac{\epsilon A^2}{8}
\end{equation}
\subsection{solution}
For the amplitude, 
\begin{equation}
    A = A_0
\end{equation}
to be determined by the initial condition. For the phase,
\begin{equation}
    \phi=\phi_0+\frac{\epsilon A^2}{8}t
\end{equation}
so the averaged solution is 
\begin{equation}
    x(t)\simeq A_0\cos\ins{\ins{1+\frac{\epsilon A^2}{8}}t+\phi_0}
\end{equation}
in leading order, and the inital conditions have 
\begin{equation}
    x(0) = A_0\cos\phi_0,\quad\dot{x}(0)=-A_0\sin\phi_0
\end{equation}
which we can combine and rearrange to
\begin{equation}
    A_0 \simeq \sqrt{x(0)^2+\dot{x}(0)^2}
\end{equation}
So the slow flow predicts that the amplitude remains constant at its initial value. There is no attracting value of \(A\) meaning the system also doesn't have a limit cycle. Looking at the frequency 
\begin{equation}
    \omega = 1+\frac{\epsilon A_0^2}{8}
\end{equation}
we see it only depends on the initial amplitude, so once that is set, the frequency is also constant.